\documentclass[11pt]{article}
\usepackage[a4paper,margin=26mm]{geometry}
\usepackage{amsmath,amssymb,amsthm,hyperref}
\newtheorem{theorem}{Theorem}
\newtheorem{lemma}{Lemma}
\newtheorem{corollary}{Corollary}
\title{Local boundary rigidity for monotone scalar transfer}
\author{Research note: a partial result, not the general transfer theorem}
\date{30 September 2026}
\begin{document}\maketitle

For a $K\times m$ real array $x=(x_j(q))$, assume each column is weakly
increasing and all distinct-coordinate mixed moments depend only on their
degree:
\[
 K^{-1}\sum_q\prod_{j\in S}x_j(q)=\mu_{|S|}.
\]
The scalar-transfer question asks whether there are $K$ real numbers $t_q$
with $K^{-1}\sum_qt_q^s=\mu_s$ for $1\le s\le m$, and, when the array lies
in an interval, whether the $t_q$ can be chosen in that interval.
The result below concerns regular scalar boundaries only.

\begin{lemma}\label{abs}
If $f,g$ are increasing functions on a probability interval, both with
mean zero, then
\[
 E|fg|\le3E(fg).
\]
In particular their covariance is nonnegative; if both are nonconstant,
it is strictly positive.
\end{lemma}
\begin{proof}
Write an increasing mean-zero function as a positive Stieltjes combination
of $h_s(u)={\bf1}_{u\ge s}-(1-s)$, $0<s<1$. For $s\le t$,
\[
 E(h_sh_t)=s(1-t),\qquad
 E|h_sh_t|=s(1-t)[1+2(t-s)]\le3s(1-t).
\]
Integrate the first equality and the inequality against the two positive
Stieltjes measures, using the triangle inequality for the absolute value.
For finite ordered probability spaces this follows by embedding the
functions as step functions. The covariance kernel is strictly positive
at every pair of interior jumps, proving the last assertion.
\end{proof}

\begin{theorem}\label{local}
Fix $m\ge2$ and a scalar array $a(q)$ with distinct levels
$a_1<\cdots<a_r$, where $r\le m-1$. Let their multiplicities be $k_l$.
Suppose a monic polynomial
\[
 P(t)=t^m+\sum_{s=0}^{m-1}\gamma_st^s
\]
satisfies $P'(a_l)=0$ at every level and $P''(a_l)>0$ at every level
whose multiplicity is at least two. Then there is a neighborhood of the
scalar array $(a(q),\ldots,a(q))$ with the following property.
Every monotone exchangeable-moment array in that neighborhood satisfying
\[
 \mu_s=K^{-1}\sum_q a(q)^s\qquad(1\le s\le m-1)
\]
satisfies
\[
 \mu_m\ge K^{-1}\sum_q a(q)^m.
\]
If equality holds, then on each scalar level all but at most one column
are constant, and every column has mean $a_l$ on that level. Such an
array has exactly the same mixed moments as the scalar array in all
degrees at most $m$.
\end{theorem}
\begin{proof}
All estimates below have constants depending only on the fixed scalar
array and $m$. On the block $B_l$ of $k_l$ rows write
\[
 x_j=a_l+b_{jl}+e_j,\qquad E_l e_j=0,
\]
where $E_l$ is the uniform mean on the block. Every $e_j$ is increasing
there. Set
\[
 C_{lij}=E_l(e_ie_j)\ge0,
 \quad S=\sum_l\sum_{i<j}C_{lij},
 \quad B=\max_{j,l}|b_{jl}|,
 \quad \delta=\max_{j,q}|x_j(q)-a(q)|.
\]
For any set $I$ of at least two indices, Lemma~\ref{abs} implies, for
any two distinct $i,j\in I$,
\[
 \left|E_l\prod_{v\in I}e_v\right|
 \le3(2\delta)^{|I|-2}C_{lij}.
\tag{1}
\]
Thus the difference between any mixed moment of $x$ and the corresponding
mixed moment of the block-constant means $a_l+b_{jl}$ is $O(S)$,
uniformly for arrays in a fixed small neighborhood.

We claim that $B=O(S)$. Consider the map from the $mr$ variables $b_{jl}$
to all mixed moments of degrees $1,\ldots,m-1$ of the block-constant
array. Its derivative at zero is injective. Indeed, a vector $b$ in its
kernel satisfies, for every $s\le m-1$ and every $s$-element set $T$,
\[
 \sum_{j\in T}\sum_l(k_l/K)a_l^{s-1}b_{jl}=0.
\]
Since $1\le s<m$, comparing sets differing in one element shows that
all the inner sums are equal, and then each is zero. For each fixed $j$,
the first $r$ of these equations form an invertible weighted Vandermonde
system, so every $b_{jl}=0$. The derivative therefore has a bounded left
inverse. The prescribed lower mixed moments and (1) now give
$B\le C(S+B^2)$; after shrinking the neighborhood this yields $B=O(S)$.

For any polynomial $P$, its symmetric multiaffine polarization in $m$
variables evaluated at $a_l+d_1,\ldots,a_l+d_m$ is
\[
 P(a_l)+\sum_{s=1}^m
       \frac{P^{(s)}(a_l)}{s!}\,
       \frac{e_s(d_1,\ldots,d_m)}{\binom ms}.
\tag{2}
\]
Here $e_s$ is the elementary symmetric polynomial. Average (2) over
rows, putting $d_j=b_{jl}+e_j$. The $s=1$ term vanishes because
$P'(a_l)=0$. The quadratic term is
\[
 \sum_l\frac{k_l}{K}\frac{P''(a_l)}{m(m-1)}
       \sum_{i<j} C_{lij}+O(S^2).
\tag{3}
\]
By (1) and $B=O(S)$, every term of degree at least three has total
absolute value $O(\delta S+S^2)$. On a singleton block all $C_{lij}$
vanish. Hence the main term in (3) is at least $cS$ for some fixed
$c>0$, unless there are no repeated levels, in which case $S=0$.
For sufficiently small $\delta$, the full difference is nonnegative,
and is strictly positive if $S>0$.

By exchangeability of mixed moments, the average polarization is
$\mu_m+\sum_{s<m}\gamma_s\mu_s$. The lower moments agree with those of
$a$, giving the desired comparison of the degree-$m$ moments. If
$S=0$, then $B=0$. Strictness in Lemma~\ref{abs} says that on each
block at most one $e_j$ is nonconstant. Every distinct-coordinate
product therefore has block mean $a_l^{|T|}$, as asserted.
\end{proof}

\begin{corollary}[Few-level rigidity]
Suppose the scalar baseline has $r$ distinct levels and $2r\le m-1$.
Every sufficiently nearby monotone exchangeable-moment array with the
same first $m-1$ moments has exactly the same mixed moments in all
degrees at most $m$. On each baseline level at most one column varies,
and all column means there equal the baseline level.
\end{corollary}
\begin{proof}
Use the proof of Theorem~\ref{local} with
$Q(t)=\prod_{l=1}^r(t-a_l)^2$ in place of $P$. Its degree is at most
$m-1$, so the expectation of its polarization is fixed by the prescribed
lower moments. Also $Q'(a_l)=0$ and
$Q''(a_l)=2\prod_{v\ne l}(a_l-a_v)^2>0$ at every level.
The same expansion gives a difference at least $cS$ for small enough
neighborhoods. Since the difference is zero, $S=0$, and the equality
argument in Theorem~\ref{local} applies.
\end{proof}

At a regular scalar constrained local minimum of the $m$-th
moment with the first $m-1$ moments fixed and $r=m-1$ distinct
interior levels, the Lagrange multiplier polynomial has precisely the
form in Theorem~\ref{local}.
Its derivative vanishes at all interior levels. Splitting two equal
nodes in opposite directions leaves the lower moments unchanged to
first order; the scalar second-order necessary condition gives
$P''(a_l)\ge0$ at every repeated level. At a regular boundary with
$r=m-1$ distinct interior levels, the critical points are simple, so
this inequality is strict. Thus the theorem rules out a local
outward crossing at every such boundary.

More precisely, when $r=m-1$, the map from the $r$ levels to the first
$r$ scalar moments has an invertible weighted Vandermonde derivative.
Nearby lower moments determine a nearby scalar branch with the same
multiplicities. Applying Theorem~\ref{local} uniformly to that branch
shows that nearby monotone tensors lie on its allowed side.

This is \emph{not} a proof of general scalar transfer. It does not
classify singular scalar boundary strata, endpoint strata, or global
components of the tensor moment image. In particular affine contraction
of all columns preserves exchangeability and can move a hypothetical
real-line counterexample arbitrarily near a fully collapsed scalar
array. The regular-boundary result alone cannot exclude that possibility.
\end{document}
